% ----------------------------------------------------------------
% ECO 701 Microeconomic Theory - Fall 2026 syllabus
% Structure and preamble follow Y:\Hunter\eco100\syllabus\eco100-syllabus.tex
% Content adapted from ECO701 syllabus 2018a (legacy/syllabus/), updated to the
% 11-topic sequence in TOPIC_PLAN.md and the Bowles/Halliday 2022 edition.
%
% Build:  latexmk -pdf eco701-syllabus.tex
% ----------------------------------------------------------------
\documentclass[11pt]{article}
\usepackage{multicol}
\usepackage{amsmath,amsthm,amssymb}
\usepackage{amsfonts}
\usepackage{graphicx}
\usepackage{tabularx}
\usepackage{enumitem}
\usepackage{xcolor}
\usepackage[colorlinks=true,linkcolor=blue,urlcolor=blue]{hyperref}
\usepackage[margin=1in]{geometry}
\usepackage{titlesec}
\usepackage[utf8]{inputenc}

\titlelabel{}
\titlespacing{\section}{0pt}{6pt}{-8pt}
\setlength{\parindent}{0pt}

\setlist[enumerate,1]{leftmargin=*, itemsep=9pt, label=\textbf{\arabic*.}, font=\bfseries, before=\bfseries}
\setlist[enumerate,2]{noitemsep, nolistsep, leftmargin=*, labelindent=-1em, label=\alph*., font=\normalfont, before=\normalfont}
\setlist[itemize]{noitemsep, nolistsep, leftmargin=*}
\setlist[description]{noitemsep, nolistsep, leftmargin=*, labelindent=-1em, font=\normalfont\scshape\color{blue}, before=\normalfont}

% Placeholder marker for details not yet confirmed by the registrar.
\newcommand{\TBD}{\textcolor{red}{\textbf{TBA}}}

\begin{document}

\begin{center}
\begin{Large}
\textbf{
Microeconomic Theory (ECO 701) \\
Hunter College \\
Fall 2026\\
}
\end{Large}
\end{center}

\textbf{Professor Jonathan Conning \hfill Meeting Time \& Place:} \\
Department of Economics \hfill Mondays\enspace 5:30 -- 7:30pm, HW 113 \\
695 Park Avenue, 1504A \hfill \textbf{Office Hours:} \\
Phone (212) 772-5403 \hfill Mondays\enspace 7:30 -- 8:30pm (HW 1504A) \\
Email: {\href{mailto:jconning@hunter.cuny.edu}{jconning@hunter.cuny.edu}} \hfill Wednesdays\enspace 2:30 -- 3:30pm, and by appt \\
\textbf{Teaching Assistant:} \TBD \hfill \textbf{Course site:} \href{https://jhconning.github.io/hunter_eco701/}{jhconning.github.io/hunter\_eco701}

\setlength{\parskip}{12pt}

\section{Course description}
This is a Master's level course in modern microeconomic theory. We adopt a critical approach to modeling the
behavior of actors who optimize subject to constraints, derive their best
response functions, and evaluate whether the resulting equilibrium interactions
lead to social outcomes that are efficient or otherwise desirable. Much of
economics, and of economic policy, can be thought of as the study of mechanism design and performance: how the terms of a contract or the rules of the game shape the behavior they govern, and how an employer, a platform, a government, or society at large, shape these so that the equilibrium interaction of individually motivated agents leads to outcomes closer to what a principal, or society, wants. Which criterion we use to judge that is itself a question we take seriously. A competitive market is one such set of rules or mechanism, remarkable in many ways, but not the only possible mechanism and not always available. When markets deliver and when they do not is a question we return to repeatedly. New technologies are producing new types of markets and social dilemmas, and we will look at how the microeconomist's toolkit is being adapted to new questions such as how humans and AI agents substitute for or complement one another, and how to design or evaluate the emerging systems of incentives in which they transact and negotiate. The applications may be new, but the questions are often old ones.

Topics include constrained optimization by the methods of Lagrange and
Kuhn-Tucker; monopoly, duopoly and competition among many firms as games with
continuous actions; production, cost and the theory of the firm; preferences,
choice and consumer demand; choice over time and under uncertainty; household
models; general equilibrium and welfare; externalities, public goods and the
commons; and the theory of asymmetric information and incentive contracts and mechanism design.

Economic intuition and understanding will always be our focus but mathematics is
an important tool in this course. Prerequisites include undergraduate
intermediate microeconomics (ECO 300 or equivalent) and enough math background
to be comfortable using and understanding integrals and derivatives of
multi-variable functions (Math 150 and 160 or 111, or equivalent); ECO 700 is a
co-requisite. The first weeks of the course are a good time to brush up on or
deepen those skills. If you are struggling, reach out early to find extra
materials and support.

You will be exposed to some applied economic modeling for simulation and
visualization using scientific Python. Jupyter notebooks accompany several of
the topics, and you are strongly encouraged to work with them.

\newpage

\section{Learning objectives}
At the completion of this course you should be able to:\vspace*{-12pt}
\begin{itemize}
\item set up and analyze economic choice problems as constrained optimization problems with one or more constraints, applying the methods of Lagrange and Kuhn--Tucker;
\item employ optimization methods to derive behavioral relationships --- individual consumer demand and expenditure functions, firms' factor demand, output supply and cost functions --- and understand their properties;
\item set up and analyze problems involving choice over time and under uncertainty;
\item analyze equilibrium models with differentiated agents: solve for market equilibrium prices and quantities, and conduct comparative static analysis of how outcomes respond to parameter changes;
\item model and predict the outcomes of strategic social interactions, employing equilibrium concepts including Nash and subgame-perfect Nash equilibrium;
\item understand the theory of incentive contracting under asymmetric information and solve for simple contracts;
\item read and understand journal articles that use the concepts and methods introduced in the course.
\end{itemize}

Completion of the course also advances the Master's program learning objectives: to ``explain core concepts in microeconomic theory,'' and, if desired, to ``be prepared to pursue doctoral study in economics or related fields such as public policy or demography.''

\section{Texts and course materials}
Two books, used in parallel. Neither needs to be purchased new.

% before=\normalfont overrides the preamble's before=\bfseries (line ~26), which
% would otherwise bold each item's whole body rather than just the title.
% If the descriptions ever come back bold, this option has been lost.
\begin{enumerate}[before=\normalfont]
\item \textbf{Bowles \& Halliday (2022)}, \emph{Microeconomics: Competition, Conflict, and Coordination}, Oxford University Press. 

A free PDF is available as an \href{https://sites.santafe.edu/~bowles/wp-content/uploads/BH_22_merged.pdf}{online PDF}, or softcover purchase. We use it for the institutional framing, game-theoretic material, and several applications. It is mathematically gentler than Nicholson/Snyder and considerably less dry.

\item \textbf{Nicholson \& Snyder}, \emph{Microeconomic Theory: Basic Principles and Extensions}, South-Western. 

Available to rent or buy in hardcover or e-book, and other formats. This is our reference for the calculus-based derivations. Any edition from the 11th forward is fine, though chapter numbers shift between editions.
\end{enumerate}

I will distribute lecture slides and handouts, along with other readings,
required and optional, that are not found in either book; we will also not cover
every topic in the texts. 

The current reading schedule, the slides, and links to additional materials are on the course site,
\href{https://jhconning.github.io/hunter_eco701/}{jhconning.github.io/hunter\_eco701}.
Look at the \href{https://jhconning.github.io/hunter_eco701/schedule}{reading
schedule} to see what is assigned next; it also lists any articles or other
material mentioned in lecture. Problem sets and submissions will use Gradescope (details to follow). Brightspace will be used for course announcements. Come to class having read the assigned material.

\section{Grading}
Your grade is based on:

\begin{center}
\begin{tabular}{lc}
\hline
\textbf{Assessment component} & \textbf{Weight}\\
\hline
Problem sets and participation & 25\%\\
Higher score of the midterm or final exam & 45\%\\
Lower score of the midterm or final exam & 30\%\\
\hline
\end{tabular}
\end{center}

Weighting the better exam more heavily is deliberate: a weak start you recover from should not define your grade.

Exam dates appear in the schedule below. \textbf{There are no makeup exams.} Absence from an exam is excused only with documented evidence of serious illness or emergency; in the case of an excused absence, performance on the other exam is re-weighted for the final grade. Final letter grades follow the official \href{https://www.hunter.cuny.edu/students/registration/records-and-transcripts/grading-structure/}{Hunter College Grading Structure}.

\section{Problem sets}
Problem sets are assigned regularly and graded on a simple basis (check plus, check, check minus, no submission). They are an important element of the course: they give you practice applying the material and prepare you for the kinds of problems you will meet on exams. You are strongly encouraged to work in study groups to discuss problems, but you must write up your own answers.

\section{Office hours, and other support}
Consult me after class with immediate questions. My office hours are Mondays 7:30--8:30pm in Hunter West 1504A, immediately following our class, and Wednesdays 2:30--3:30pm, as well as by appointment. The Teaching Assistant will hold separate office hours. Schedules and updates will be posted to Brightspace.

\newpage

\section{Eco 701 Course outline and reading schedule}
The class meets \textbf{Mondays, 5:30 -- 7:30pm in HW 113}. I reserve the right to adjust pace and coverage. Always check the latest
online \href{https://jhconning.github.io/hunter_eco701/schedule}{reading
schedule} each week.

\begin{center}
\begin{tabularx}{\textwidth}{@{}c l X l@{}}
\hline
\textbf{\#} & \textbf{Date} & \textbf{Topic} & \textbf{Reading}\\
\hline
1  & Mon Aug 31 & Course logistics. Definitions and key tools: institutions, mechanisms, games, equilibrium, efficiency and distribution & BH 1\\
2  & Mon Sep 14 & Definitions and key tools, concluded: best response and Nash equilibrium. \textbf{Preferences and choice}: rationality axioms, utility representation and its limits, ordinality, monotonicity and convexity, the indifference-curve zoo & BH 1--2\\
3  & Mon Sep 28 & Constrained optimization I: the Lagrange method, comparative statics, the envelope theorem & NS 2\\
4  & Mon Oct 5  & Constrained optimization II: inequality constraints and Kuhn--Tucker & NS 2--3\\
5  & \emph{Tue} Oct 13 & Production optimization: cost minimization, profit maximization, duality & NS 9--11\\
6  & Mon Oct 19 & Consumption optimization: utility maximization, demand and expenditure functions & NS 4--5\\
7  & Mon Oct 26 & Consumption optimization, continued: income and substitution effects, labor supply & NS 5\\
8  & Mon Nov 2  & \textbf{Midterm exam} (tentative) & ---\\
9  & Mon Nov 9  & Choice over time and under uncertainty: discounting, expected utility, risk sharing & NS 7, 17\\
10 & Mon Nov 16 & Household economics: the farm-household model, separability and its failure & Handout\\
11 & Mon Nov 23 & General equilibrium I: exchange, the core, the welfare theorems & NS 13\\
12 & Mon Nov 30 & General equilibrium II: production and two-sector models & NS 13, BH\\
13 & Mon Dec 7  & Externalities, public goods, and the commons; $N$-player games & NS 19, BH\\
14 & Mon Dec 14 & Asymmetric information: moral hazard and the principal--agent problem. How micro theory frames empirical work. Final review & NS 18, BH 5, 10\\
\hline
   & Mon Dec 21 & \textbf{Final exam}, 5:30 -- 7:30pm, HW 113 & ---\\

\hline
\end{tabularx}
\end{center}

\newpage

\section{Use of artificial intelligence}
AI assistants (Claude, ChatGPT, Gemini, Deepseek, and others) can be genuinely
useful: for explaining derivations you might be unsure of, checking your algebra, creating extra practice problems, or quizes before an exam. They can help you stay organized, and lead you to explore more advanced material or the history of ideas. Used well, they can be a strong complement to your learning efforts, and I encourage it.  

The problem is that the same tools make it easy to produce work you do not understand. There is now \href{https://www.economist.com/graphic-detail/2026/08/18/does-ai-stop-children-from-learning}{strong quasi-experimental evidence} that students who take shortcuts on homework using AI, later score noticeably lower on exams  (and economists are already building models of possible societal "\href{https://www.nber.org/papers/w34910}{knowledge collapse"}). Do not outsource your thinking to AI. Two rules follow:

\begin{itemize}
\item \textbf{No AI during exams.} Exams are closed-book and closed-device.
\item \textbf{Problem sets are your own write-up.} Consult as you like while
working, including AI and classmates, but hand in your own
reasoning in your own words. I read for the intermediate steps, the graph you
drew, the sentence explaining why the constraint binds, etc. Polished answers with no evidence of work earns little credit, and tells me nothing about what you can do.
\end{itemize}

If you use an AI tool substantially on a problem set, note it briefly at the top of your submission. This costs you nothing and keeps everyone honest.  Using AI on exams or submitting AI work as your own without acknowledgement is academic dishonesty and will be treated according to the College's policy. 

The content of this course has evolved over many years of teaching related material.  As of this semester, I am using AI (tools such as Claude Code and Google Antigravity, and Mathpix) to stay organized and generate and post more polished content to the class site. I have used it to generate or improve some problems and examples, and illustrative plots and simulations using python. Much content and scaffolding is posted at \href{https://github.com/jhconning/hunter_eco701}{https://github.com/jhconning/hunter\_eco701}.



\section{Academic integrity}
The following is \href{https://www.hunter.cuny.edu/students/student-affairs/office-of-student-conduct/academic-integrity/}{Hunter College policy}:

Academic dishonesty will not be tolerated. No dictionaries, cell phones, headphones, or electronic devices of any kind are allowed during exams, nor is talking to other students, however briefly or for whatever reason.

Hunter College regards acts of academic dishonesty (e.g., plagiarism, cheating on examinations, obtaining unfair advantage, and falsification of records and official documents) as serious offenses against the values of intellectual honesty. The College is committed to enforcing the CUNY Policy of Academic Integrity and will pursue cases of academic dishonesty according to the Hunter College Academic Integrity Procedures.

\section{ADA policy}
In compliance with the American Disability Act of 1990 (ADA) and with Section 504 of the Rehabilitation Act of 1973, Hunter College is committed to ensuring educational parity and accommodations for all students with documented disabilities and/or medical conditions. It is recommended that all students with documented disabilities (Emotional, Medical, Physical, and/or Learning) consult the Office of AccessABILITY, located in Room E1214B, to secure necessary academic accommodations. For further information and assistance, please call: (212) 772-4857 or (212) 650-3230.

\section{Policy on sexual misconduct}
In compliance with the CUNY Policy on Sexual Misconduct, Hunter College reaffirms the prohibition of any sexual misconduct, which includes sexual violence, sexual harassment, and gender-based harassment retaliation against students, employees, or visitors, as well as certain intimate relationships. Students who have experienced any form of sexual violence on or off campus (including CUNY-sponsored trips and events) are entitled to the rights outlined in the Bill of Rights for Hunter College.

Sexual Violence: Students are strongly encouraged to immediately report the incident by calling 911, contacting the NYPD Special Victims Division Hotline (646-610-7272) or their local police precinct, or contacting the College's Public Safety Office (212-772-4444).

All Other Forms of Sexual Misconduct: Students are also encouraged to contact the College's Title IX Campus Coordinator, Dean John Rose (\href{mailto:jtrose@hunter.cuny.edu}{jtrose@hunter.cuny.edu} or 212-650-3262) or Colleen Barry (\href{mailto:colleen.barry@hunter.cuny.edu}{colleen.barry@hunter.cuny.edu} or 212-772-4534) and seek complementary services through the Counseling and Wellness Services Office, Hunter East 1123.

\href{http://www.cuny.edu/about/administration/offices/la/Policy-on-Sexual-Misconduct-12-1-14-with-links.pdf}{CUNY Policy on Sexual Misconduct}

\end{document}
